% answers_short.tex —— 2014 年·天津卷·理科 · 简短答案速查
% 由 tools/make_short_answers.py 自 answers.tex 的 \gksolhead/\gkanswer 提取，
% 并与 2014/2014_数学真题_天津理_答案要点.md 逐题比对后生成，勿手改（改源 MD 后重跑）。

\begin{gknotice}
\textbf{使用说明}：本节只给最终答案便于快速核对，完整推导见解析版。答案为编写者独立推导，\textbf{非官方标准答案}。
\end{gknotice}

\gkpart{选择题}
\begin{center}\begin{tabular}{ccccccccc}
\toprule
题号 & 1 & 2 & 3 & 4 & 5 & 6 & 7 & 8 \\
\midrule
答案 & A & B & B & D & A & D & C & C \\
\bottomrule
\end{tabular}\end{center}

\gkpart{填空题}
\begin{itemize}[leftmargin=2.6em]
  \item[\textbf{9.}] $60$。总份数 $20$，每份 $15$ 人，一年级 $4$ 份即 $60$ 人。
  \item[\textbf{10.}] $\dfrac{20\pi}{3}$。几何体为半径 $2$、高 $2$ 的圆锥与半径 $1$、高 $4$ 的圆柱同轴叠放，$\dfrac{8\pi}{3}+4\pi=\dfrac{20\pi}{3}$。
  \item[\textbf{11.}] $-\dfrac12$。由 $S_2^2=S_1S_4$ 得 $(2a_1-1)^2=a_1(4a_1-6)$，解得 $a_1=-\dfrac12$。
  \item[\textbf{12.}] $-\dfrac14$。由 $2\sin B=3\sin C$ 得 $b=\dfrac32c$，再由 $b-c=\dfrac a4$ 得 $a=2c$，用余弦定理得 $\cos A=-\dfrac14$。
  \item[\textbf{13.}] $3$。圆为 $x^2+(y-2)^2=4$，弦 $y=a$ 的半弦长 $\sqrt{4-(a-2)^2}$；由 $|AB|=|OA|$ 得 $4-(a-2)^2=a$，即 $a=3$。
  \item[\textbf{14.}] $(0,1)\cup(9,+\infty)$。在 $(-3,0)$ 内根的个数为 $0<a<1$ 时 $2$ 个、$a=1$ 时 $1$ 个、$a>1$ 时 $0$ 个；在 $(1,+\infty)$ 内为 $a>9$ 时 $2$ 个、$a=9$ 时 $1$ 个、$a<9$ 时 $0$ 个；另两支各恒有 $1$ 个根，合计恰为 $4$ 个当且仅当 $a\in(0,1)\cup(9,+\infty)$。
\end{itemize}

\gkpart{解答题}
\begin{itemize}[leftmargin=2.6em]
  \item[\textbf{15.}] (1) $\pi$。化简得 $f(x)=\dfrac12\sin\left(2x-\dfrac\pi3\right)$，故最小正周期为 $\dfrac{2\pi}{2}=\pi$；(2) 最大值 $\dfrac14$，最小值 $-\dfrac12$。当 $2x-\dfrac\pi3\in\left[-\dfrac{5\pi}{6},\dfrac\pi6\right]$ 时，$\sin$ 的最大值为 $\dfrac12$（$x=\dfrac\pi4$），最小值为 $-1$（$x=-\dfrac{\pi}{12}$）
  \item[\textbf{16.}] (1) $\dfrac{49}{60}$。数学学院的 $3$ 人同属一个学院，其余 $7$ 人各属不同学院；符合条件的结果数为 $\mathrm C_7^3+3\mathrm C_7^2=98$，而总数为 $\mathrm C_{10}^3=120$；(2) $P(X=0)=\dfrac16$，$P(X=1)=\dfrac12$，$P(X=2)=\dfrac3{10}$，$P(X=3)=\dfrac1{30}$，$E(X)=\dfrac65$。按超几何分布计算，或直接取 $E(X)=3\times\dfrac4{10}=\dfrac65$
  \item[\textbf{17.}] (1) 见解析版；(2) $\dfrac{\sqrt3}{3}$。平面 $PBD$ 的法向量取 $(2,1,1)$，$\sin\theta=\dfrac{|\overrightarrow{BE}\cdot\boldsymbol n|}{|\overrightarrow{BE}||\boldsymbol n|}=\dfrac{4}{\sqrt2\cdot2\sqrt6}=\dfrac{\sqrt3}{3}$；(3) $\dfrac{3\sqrt{10}}{10}$。由 $\overrightarrow{BF}\cdot\overrightarrow{AC}=0$ 得 $F$ 分 $PC$ 的比为 $\dfrac14$，两半平面法向量分别取 $(0,1,0)$ 与 $(0,-3,1)$，其余弦绝对值为 $\dfrac{3}{\sqrt{10}}$，二面角为锐角
  \item[\textbf{18.}] (1) $\dfrac{\sqrt2}{2}$。由 $a^2+b^2=3c^2$ 与 $b^2=a^2-c^2$ 得 $a^2=2c^2$，故 $e=\dfrac{\sqrt2}{2}$；(2) $4+\sqrt{15}$ 或 $4-\sqrt{15}$。由 $\overrightarrow{F_1P}\cdot\overrightarrow{F_1B}=0$ 得 $P\left(-\dfrac{4c}{3},\dfrac c3\right)$，圆心为 $\left(-\dfrac{2c}{3},\dfrac{2c}{3}\right)$、半径为 $\dfrac{\sqrt5c}{3}$，令 $l:y=kx$ 用点到直线距离等于半径，得 $k^2-8k+1=0$
  \item[\textbf{19.}] (1) $A=\{0,1,2,3,4,5,6,7\}$。$q=2$、$n=3$ 时 $x=x_1+2x_2+4x_3$，三系数各取 $0$ 或 $1$，共 $2^3=8$ 个互异值；(2) 见解析版
  \item[\textbf{20.}] (1) $\left(0,\dfrac1{\mathrm e}\right)$。$f'(x)=1-a\mathrm e^x$，$a>0$ 时最大值 $f(-\ln a)=-\ln a-1$，有两个零点须 $-\ln a-1>0$；(2) 见解析版；(3) 见解析版
\end{itemize}
